\chapter{Conclusión}

El sistema desarrollado responde directamente a los tres síntomas identificados en el
caso BookPoint Chile: la lentitud de procesos se aborda mediante microservicios
independientes que escalan y se despliegan sin afectarse entre sí; el stock
desactualizado entre sucursales se resuelve centralizando el inventario en un servicio
único consultado por todos los canales; y las demoras de despacho se abordan con un
microservicio de envíos dedicado, capaz de registrar y exponer el historial completo de
estados de cada envío.

La arquitectura resultante -- diez microservicios independientes, cada uno con su
propia base de datos, comunicándose vía Feign y descubriéndose mediante Eureka detrás de
un API Gateway -- cubre en profundidad el flujo comercial completo del caso: explorar el
catálogo, agregar productos al carrito, comprar, pagar, despachar y reseñar. Quedan
deliberadamente fuera de alcance funciones administrativas de soporte (autenticación y
roles, reportes de venta, promociones, transferencias entre sucursales), priorizando
profundidad y calidad en el flujo principal por sobre cobertura superficial de la
totalidad del caso.

Un aspecto que vale la pena destacar de este proceso es que las correcciones
documentadas en la sección~\ref{sec:evolucion} no fueron anticipadas desde el diseño
inicial, sino descubiertas y corregidas iterativamente a medida que se sometía el
sistema a verificación real -- probando cada flujo contra el stack completo levantado
en Docker, en vez de asumir que el código funcionaba por haber compilado. Esa disciplina
de verificación empírica es, en sí misma, una de las buenas prácticas que este informe
buscó demostrar.

% Espacio para agregar una reflexión personal sobre aprendizajes del proceso, si el
% estudiante quiere incorporar algo aquí antes de la entrega final.
